% Lab 9 — SI Noise Analysis（原书 PDF p.103–114）
\bichapter{Lab 9：SI 噪声分析}{Lab 9: SI Noise Analysis}

\begin{objectiveBox}
完成本实验后，你应能够：
\begin{itemize}
  \item 运行 PrimeTime SI 进行串扰噪声（crosstalk noise）分析；
  \item 在 shell 与 GUI 中生成并解读关键噪声报告。
\end{itemize}
\end{objectiveBox}

\begin{durationBox}
约 \textbf{45} 分钟。
\end{durationBox}

\section{概述}
\label{lab9:overview}

\noindent\textit{Overview}

本实验涵盖：噪声分析运行脚本要点、shell 中的噪声报告解读，以及 GUI 中的噪声 bump 直方图分析。

\subsection{相关文件与目录}
\label{lab9:files}

\noindent\textit{Relevant Files and Directories}

\begin{tabularx}{\textwidth}{@{}l X@{}}
\toprule
\textbf{路径 / 文件} & \textbf{说明} \\
\midrule
\cmd{ref/db/} & 工艺库 \\
\cmd{ref/design_data/} & 寄生参数与设计网表 \\
\cmd{ref/scripts/} & 支持脚本 \\
\cmd{ref/tcl_procs/} & 实用 Tcl 过程 \\
\cmd{lab9_si_noise/} & 本实验工作目录 \\
\quad\cmd{run_noise.pt} & PTSI 噪声运行脚本 \\
\quad\cmd{.synopsys_pt.setup} & Setup 文件 \\
\quad\cmd{SE_PE_savesession/} & 含噪声分析的已保存会话 \\
\quad\cmd{noise.log} & SI 运行日志 \\
\quad\cmd{rpt/} & 运行脚本生成的报告 \\
\bottomrule
\end{tabularx}


\section{操作说明}
\label{lab9:instructions}

\noindent\textit{Instructions}

\subsection{任务 1：运行 PrimeTime SI 噪声分析}
\label{lab9:task1}

\noindent\textit{Task 1. Run PrimeTime SI for Noise Analysis}

噪声运行脚本已预先执行并记录日志。

\begin{enumerate}
  \item 进入目录 \cmd{lab9_si_noise}。

  \item 打开运行脚本 \cmd{run_noise.pt} 并回答：
\begin{questionBox}
\textbf{问题 1.}\ 列出该脚本中与噪声分析相关的重要步骤（至少三条）。\\[0.35em]
\textbf{答：} 与噪声分析相关的步骤示例（任选其三即可）：
\begin{itemize}
  \item 用变量 \cmd{si\_enable\_analysis} 开启 SI 分析（与 SI 延迟分析共用）；
  \item 使用包含噪声表征的库；
  \item 应用噪声 derating 因子缩放计算的噪声 bump 高度与宽度；
  \item 读入网寄生参数并保留耦合电容（与 SI 延迟分析共用）；
  \item 执行 \cmd{update\_noise} 触发噪声更新（或执行 \cmd{report\_noise}，其会同时更新并报告）。
\end{itemize}
\end{questionBox}

\begin{questionBox}
\textbf{问题 2.}\ 勾选噪声分析默认开启的选项（由 \cmd{set\_noise\_parameters} 控制）：\\
\quad$\square$ Beyond rail analysis \quad$\square$ Between-the-rails analysis\\
\quad$\square$ Noise propagation \quad$\square$ Timing windows are used to filter aggressors\\[0.35em]
\textbf{答：} 默认开启：\\
\quad$\square$ Beyond rail analysis \quad$\checkmark$ Between-the-rails analysis\\
\quad$\square$ Noise propagation \quad$\checkmark$ Timing windows are used to filter aggressors
\end{questionBox}

  \item 打开日志文件 \cmd{noise.log} 并回答：
\begin{questionBox}
\textbf{问题 3.}\ 哪条命令触发了 timing update（查找消息 \textbf{UITE-214}）？\\[0.35em]
\textbf{答：} \cmd{check\_timing} 触发了 timing update。

\textbf{问题 4.}\ 该 timing update 是否包含串扰延迟分析？\\[0.35em]
\textbf{答：} \textbf{是}。可从 \textbf{XTALK-xxx} 消息看出。

\textbf{问题 5.}\ 哪条命令触发了 noise update（查找消息 \textbf{NOISE-010}）？\\[0.35em]
\textbf{答：} \cmd{update\_noise} 触发了 noise update。\\
相关消息由应用变量 \cmd{si\_noise\_update\_status\_level} 开启。

\textbf{问题 6.}\ 根据课堂讲义，timing update 为噪声分析提供哪些信息？\\[0.35em]
\textbf{答：} 使用 aggressor 网的\textbf{到达窗口（arrival windows）}进行滤波；
仅窗口重叠的 aggressor 被定义为 active，并参与噪声 bump 计算。
\end{questionBox}
\end{enumerate}

\subsection{任务 2：在 Shell 中解读噪声报告}
\label{lab9:task2}

\noindent\textit{Task 2. Interpret Reports for Noise in the Shell}

\begin{enumerate}
  \item 从 \cmd{lab9_si_noise} 目录启动 PrimeTime，并恢复 \cmd{SE_PE_savesession} 中的会话（善用 Tab 补全）：
\begin{lstlisting}
unix% pt_shell
pt_shell> restore_session SE_PE_savesession
\end{lstlisting}

  \item 执行下列命令，用 slack 类型 \texttt{area} 找出最坏噪声 slack：
\begin{lstlisting}
pt_shell> report_noise -all_violators -slack_type area
pt_shell> report_noise -verbose -slack_type area
\end{lstlisting}

\begin{questionBox}
\textbf{问题 7.}\ 本设计是否存在噪声违例？\\[0.35em]
\textbf{答：} \textbf{没有}噪声违例。

\textbf{问题 8.}\ 最坏噪声 slack（面积单位）是多少？相关区域是 \texttt{above\_low} 还是 \texttt{below\_high}？\\[0.35em]
\textbf{答：} 最坏 slack 为 \textbf{153.8006}（\texttt{area} 单位），区域为 \textbf{above\_low}。

\textbf{问题 9.}\ 对该“最坏”victim 网列出了多少个“有效（effective）”aggressor 网？\\[0.35em]
\textbf{答：} 有 \textbf{1} 个有效 aggressor 网：\cmd{n25422}。

\textbf{问题 10.}\ 对该 victim 网，给出驱动单元 pin 名、负载 pin 名与网名。\\[0.35em]
\textbf{答：} victim 网名 \cmd{n25449}；驱动 pin \cmd{T3\_reg\_190/QN}；负载 pin \cmd{U11405/A1}。

\textbf{问题 11.}\ 解释报告中 ``Propagated'' 的含义，并说明该行为何可能为空（显示为破折号）。\\[0.35em]
\textbf{答：} ``Propagated'' 行用于\textbf{前级网噪声经驱动单元传播}到本网的噪声。
该行可能为空，因为未开启噪声传播（本分析即如此），和/或库不支持噪声传播（本实验库亦不支持）。
\end{questionBox}

  \item 对同一网执行下列命令，根据报告\textbf{上部}（噪声计算概要）回答问题：
\begin{lstlisting}
# 从 report_noise 输出复制粘贴 pin 名以免拼写错误
pt_shell> report_noise_calculation -above_low -from <pin> -to <pin>
\end{lstlisting}

\begin{questionBox}
\textbf{问题 12.}\ 解释 ``Steady state resistance source: library set iv curve'' 的含义。\\[0.35em]
\textbf{答：} 表示计算噪声 bump 高度时，驱动输出稳态 I-V 特性所用的模型；\\
其他选项例如 \cmd{set\_steady\_state\_resistance} 设置的 ``user set value''，
或 PrimeTime 估计线性电阻时的 ``estimation set value''。

\textbf{问题 13.}\ 是否对该网应用了噪声 derating 因子？是否与运行脚本一致？\\[0.35em]
\textbf{答：} \textbf{是}。高度与宽度缩放因子均为 \textbf{1.5}，与运行脚本一致。

\textbf{问题 14.}\ 解释 ``noise effort threshold'' 的含义。\\[0.35em]
\textbf{答：} 噪声 bump 高度大于轨到轨电压的 \textbf{0.2} 倍（亦为默认值）的网，
将使用基于 Arnoldi 阶约化的更复杂计算器。
\end{questionBox}

  \item 根据报告\textbf{中部}（噪声计算细节）回答：
\begin{questionBox}
\textbf{问题 15.}\ 哪条命令可排除特定 aggressor 网，或将其定义为具有无限 timing window？\\
（提示：参阅 Job Aid。）\\[0.35em]
\cmd{set\_si\_noise\_analysis}。
\end{questionBox}

  \item 根据报告\textbf{下部}（噪声 slack 计算）回答：
\begin{questionBox}
\textbf{问题 16.}\ 解释 ``Constraint type: library immunity table'' 的含义。\\[0.35em]
\textbf{答：} 表示计算噪声约束（允许/所需噪声 bump 高度）所用的模型；\\
其他选项例如 \cmd{set\_noise\_margin} 的 ``user margin''，或库与用户均未提供约束时的 ``none''。

\textbf{问题 17.}\ 噪声约束（库电压单位下的所需噪声 bump 高度）是多少？\\[0.35em]
\textbf{答：} 噪声约束为 \textbf{0.6009}（库电压单位）。

\textbf{问题 18.}\ slack 的 ``area'' 与上一节噪声 bump 的 ``area'' 有何不同？\\[0.35em]
\textbf{答：} 噪声 bump area 为 \textbf{\%height $\times$ width}；\\
slack area 为 \textbf{height $\times$ width}。可用 shell 中的 \cmd{expr} 验证。
\end{questionBox}

  \item 查看所有寄存器数据 pin 上 slack 类型为 \texttt{area\_percent} 的最差 5 条网：
\begin{lstlisting}
pt_shell> report_noise -data_pins -nworst 5 -slack_type area_percent
\end{lstlisting}

\begin{questionBox}
\textbf{问题 19.}\ 最差的触发器数据 pin 距离约束有多远（裕量很大还是很小）？\\[0.35em]
\textbf{答：} 裕量较大：\\
above\_low slack：\textbf{58.27\%}（pin：\cmd{T1\_reg\_120/SD}）；\\
below\_high slack：\textbf{INFINITY}（pin：\cmd{T1\_reg\_120/SD}）。

\textbf{问题 20.}\ 解释 slack 为 \texttt{INFINITY} 的含义。\\[0.35em]
\textbf{答：} 若计算的噪声 bump 为零，则 slack 定义为 \textbf{INFINITY}。\\
可能原因：该网无有效 aggressor，或所有有效 aggressor 均被筛除。
\end{questionBox}

  \item 查看本设计中寄存器时钟 pin 上的噪声：
\begin{lstlisting}
pt_shell> report_noise -clock_pins
\end{lstlisting}
\end{enumerate}

\subsection{任务 3：在 GUI 中运行噪声分析}
\label{lab9:task3}

\noindent\textit{Task 3. Run Noise Analysis in the GUI}

独立于 slack 与约束高度，识别具有较大 bump 高度的网可能很有意义；GUI 便于此类分析。

\begin{enumerate}
  \item 打开 GUI：
\begin{lstlisting}
pt_shell> start_gui
\end{lstlisting}

  \item 通过工具栏按钮或 \textbf{Noise} 下拉菜单打开 \textbf{Accumulated Noise Bump Histogram}，
        在对话框中保留默认设置并点击 \textbf{OK}。

\begin{noteBox}
该直方图按 bump \textbf{高度}（非 slack）列出网或负载 pin（在对话框中选择），
针对 \texttt{below\_high} 区域（在对话框中选择）。\\
对本设计，\textbf{Composite Noise Bump Histogram} 结果相同；
composite 噪声 bump 包含 propagated noise，而 accumulated 噪声 bump 不包含。
\end{noteBox}

  \item 选择最坏 bin（最右侧柱）以显示 bump 高度最大的网名。

\begin{questionBox}
\textbf{问题 21.}\ 有多少条网的噪声 bump 高度大于 0.5？\\[0.35em]
\textbf{答：} 在 \texttt{below\_high} 区域，有 \textbf{4} 条网的噪声 bump 高度大于 0.5：
\begin{lstlisting}
pt_shell> report_noise_bumps -bump_greater_than 0.5
Pin Name    Net Name
U11361/A1   n25502
U13544/A2   C230
U11688/B1   AWO_199
U15062/A2   BS1_RDY
\end{lstlisting}
\end{questionBox}

\begin{noteBox}
SolvNet（Doc Id \textbf{009418}）提供可在 shell 中报告相同信息的 Tcl 过程；
实验环境已包含该过程。示例：
\begin{lstlisting}
pt_shell> report_noise_bumps -bump_greater_than 0.5
\end{lstlisting}
\end{noteBox}

  \item 退出 PrimeTime。
\end{enumerate}

\noindent 恭喜——Lab~9 完成！
